\section{Phần 2 -- Xây dựng bản đồ bằng SLAM Toolbox}

\subsection{Cách SLAM Toolbox hoạt động trong bài}
SLAM Toolbox nhận \code{/scan} và TF \code{odom}$\rightarrow$\code{base\_link}. Với mỗi scan đủ điều kiện (robot đã di chuyển
quá ngưỡng), nó \emph{scan matching} với các scan trước để hiệu chỉnh vị trí, thêm một nút vào đồ thị tư thế, rồi giải tối ưu đồ thị
(Ceres) khi có \emph{loop closure}. Kết quả là bản đồ lưới chiếm dụng trên \code{/map} và phép biến đổi \code{map}$\rightarrow$\code{odom}.
Launch dùng là \code{turtlebot4\_navigation/slam.launch.py}, chỉ đổi file cấu hình và chọn node bất đồng bộ:
\begin{lstlisting}[language=bash,caption={\file{scripts/start\_slam.sh}.}]
exec ros2 launch turtlebot4_navigation slam.launch.py use_sim_time:=true sync:=false \
  params:="$DIR/config/slam_small.yaml" "$@"
\end{lstlisting}

\subsection{Chỉnh tham số cho mô phỏng chậm}
Với cấu hình mặc định, node đồng bộ (\code{sync:=true}) xử lý từng scan nên hàng đợi tăng
(\emph{Queue size has grown to: 40}), bản đồ cập nhật chậm và nhiễu. Em dùng node bất đồng bộ (\code{sync:=false}: bỏ qua scan khi đang bận)
và chỉnh \file{config/slam\_small.yaml} từ bản gốc \code{slam.yaml}:
\begin{table}[H]\centering\small
\begin{tabular}{lcc}\toprule
Tham số & Gốc & Dùng\\\midrule
\code{throttle\_scans} & 1 & 2\\
\code{map\_update\_interval} & 0.5 & 1.0\\
\code{max\_laser\_range} & 12.0 & 10.0\\
\code{minimum\_time\_interval} & 0.25 & 0.4\\
\code{transform\_timeout} & 0.2 & 0.5\\
\code{minimum\_travel\_distance} & 0.0 & 0.05\\
\code{minimum\_travel\_heading} & 0.0 & 0.1\\
\code{scan\_buffer\_maximum\_scan\_distance} & 12.0 & 10.0\\
\bottomrule\end{tabular}
\caption{Tham số SLAM Toolbox đã thay đổi so với \code{slam.yaml} gốc (các tham số còn lại giữ nguyên: Ceres, \code{resolution: 0.05}, \code{do\_loop\_closing: true}).}
\end{table}
\code{minimum\_travel\_*} lớn hơn 0 giúp không thêm scan khi robot đứng yên (tránh nhiễu tích lũy); \code{transform\_timeout} lớn hơn để
chịu được TF đến muộn trên máy chậm.

\subsection{Điều khiển bằng teleoperation}
Bản đồ được quét bằng \code{teleop\_twist\_keyboard} với tốc độ rất thấp để scan matching ổn định:
\begin{lstlisting}[language=bash,caption={\file{scripts/teleop\_slow.sh} (dùng 0,05\,m/s và 0,2\,rad/s).}]
ros2 run teleop_twist_keyboard teleop_twist_keyboard --ros-args -p speed:="$SPEED" -p turn:="$TURN"
\end{lstlisting}
Phím: \code{i} tiến, \code{,} lùi, \code{j}/\code{l} quay trái/phải, \code{k} dừng. Lái đi dọc các tường, vòng qua hai hộp vật cản và qua
khe giữa hai khu để robot quét hết khu vực có thể di chuyển, đồng thời quan sát bản đồ hình thành trong RViz2.
Em đã thử viết thêm một node tự động bám tường nhưng \textbf{không dùng} (không cho bản đồ tốt); bản đồ nộp là do lái tay.

\anh[0.85]{fig_slam_progress.png}{Bản đồ đang hình thành trong RViz2 khi chạy SLAM Toolbox và lái robot bằng teleop: vùng xám chưa biết, vùng trắng đã quét, tia LiDAR màu đỏ trùng với tường.}

\subsection{Lưu bản đồ}
\begin{lstlisting}[language=bash,caption={\file{scripts/save\_map.sh} (rút gọn).}]
ros2 run nav2_map_server map_saver_cli -f "$DIR/maps/map" --ros-args \
  -p save_map_timeout:=20.0 -p free_thresh_default:=0.25 -p occupied_thresh_default:=0.65
\end{lstlisting}
Kết quả là hai file \file{maps/map.pgm} (ảnh $262\times 220$ điểm, 0,05\,m/điểm) và \file{maps/map.yaml}:
\begin{lstlisting}[language=yaml,caption={\file{maps/map.yaml}.}]
image: map.pgm
mode: trinary
resolution: 0.05
origin: [-4.9, -2.32, 0]
negate: 0
occupied_thresh: 0.65
free_thresh: 0.25
\end{lstlisting}
Ý nghĩa: \code{resolution} là kích thước một điểm ảnh (m); \code{origin} là tọa độ góc dưới-trái của ảnh trong frame \code{map};
điểm có xác suất chiếm dụng $>0{,}65$ là tường (đen), $<0{,}25$ là trống (trắng), còn lại là chưa biết (xám). Gốc
\code{map} trùng vị trí robot ngay sau khi undock, vì SLAM được bật ở đó.

\subsection{Hậu xử lý bản đồ (cần nêu rõ)}
Bản đồ thô từ SLAM có nhiều điểm lẻ \emph{bên ngoài} tường: tia LiDAR lọt qua khe hở nhỏ giữa các tấm tường ở góc và để lại các
vệt trống/vật cản rời rạc. Đây không phải tường nhân đôi, nhưng làm bản đồ xấu và làm Nav2 thấy ``vật cản ma''.
Em dùng \file{scripts/clean\_map.py}: (i) co (erode) vùng trống bán kính 2 điểm rồi lấy thành phần liên thông lớn nhất, (ii) nở (dilate)
trở lại trong vùng trống gốc, (iii) chỉ giữ ô tường sát vùng đó ($\le 2$ điểm), phần còn lại đặt thành ``chưa biết''.
\begin{lstlisting}[language=Python,caption={Lõi của \file{scripts/clean\_map.py}.}]
core = largest_component(erode(free, 2))
room = dilate(core, 2) & free                 # vung free that su cua arena
keep_occ = occ & dilate(room, 2)              # tuong sat vung free
out = np.full_like(img, UNK)                  # UNK = 205 (xam)
out[room] = FREE                              # FREE = 254
out[keep_occ] = OCC                           # OCC = 0
\end{lstlisting}
Kết quả: ô trống 25\,412 $\to$ 23\,048, ô tường 2\,892 $\to$ 1\,929. Cấu trúc bên trong (tường, vách, hai hộp) được giữ nguyên;
bản gốc lưu ở \file{maps/map\_raw.pgm}.
\anh{fig_map_raw_vs_clean.png}{Bản đồ thô từ SLAM (trái) và sau hậu xử lý (phải).}
\anh[0.8]{fig_map_final.png}{Bản đồ cuối cùng \file{maps/map.pgm}. Trục tính bằng mét trong frame \code{map}; hai hộp vật cản hiện là khối xám vì LiDAR chỉ thấy mặt ngoài.}

\subsection{Đánh giá}
Bản đồ có đủ tường biên, vách dọc, vách ngang (hành lang phía bắc) và hai hộp vật cản; không có tường nhân đôi hay vùng lệch rõ rệt.
Toàn bộ bản đồ bị xoay nhẹ khoảng $5^\circ$ so với hệ trục lưới, vì gốc \code{map} lấy theo hướng robot lúc bắt đầu quét, không phải theo
trục của world. Điều này không ảnh hưởng đến localization và navigation vì chúng chỉ dùng chính bản đồ này.
